\section{Project GR00T：通用机器人基础模型}

\subsection{愿景与架构}

GR00T 的目标是构建能操控各种机器人身体的通用基础模型：

\begin{itemize}
    \item 接受自然语言指令和视觉输入
    \item 输出机器人动作序列
    \item 跨机器人形态泛化——同一个模型控制不同的机器人
\end{itemize}

\subsection{核心技术挑战}

\begin{warningbox}{机器人基础模型的独特挑战}
\begin{itemize}
    \item \textbf{数据稀缺}：不像语言有互联网规模的数据，机器人操作数据极度稀缺
    \item \textbf{sim-to-real gap}：仿真中训练的策略在真实世界中往往失效
    \item \textbf{安全性}：机器人直接作用于物理世界，错误可能造成物理伤害
    \item \textbf{实时性}：物理交互要求毫秒级的决策延迟
\end{itemize}
\end{warningbox}

\subsection{数据生成策略}

\begin{itemize}
    \item 利用大规模物理仿真（Isaac Sim）生成训练数据
    \item 通过 LLM 自动生成多样化的任务和 reward function（Eureka 方法）
    \item 从互联网视频中学习人类操作（观察学习）
    \item Human demonstration + teleoperation 收集真实操作数据
\end{itemize}

\subsection{本章小结}

GR00T 代表了将基础模型范式从语言扩展到物理世界的雄心，核心挑战在于数据、sim-to-real 迁移和安全性。

